\documentclass{book}
\usepackage{hyperref}
\usepackage[margin=1in]{geometry}
\usepackage{ulem}
\usepackage{tipa} % Add this line to support \textipa command
\hypersetup{
    colorlinks=true,
    linkcolor=black,
    filecolor=magenta,      
    urlcolor=cyan,
    }

\urlstyle{same}

\title{French Conjugation}
\author{Junzhang Luo}

\begin{document}

\tableofcontents

\chapter{Day 1}
    \section{Introduction}
        Bonjour, je m'appelle Jz. Je suis un étudiant en informatique de l'université de Waterloo au Canada.
        J'ai dix-huit ans. J'habite à l'appartement à côte du campus de l'université. 
        J'aime faire du sports, la musique et j'adore explorer la technologie de pointe, c'est l'intelligence artifacielle et les applications.
        Je parle anglais, chinois couramment et je suis en train de apprendre français. Mon objectif est de passer l'examen DELF nievu B2 ou C1 
        à la fin de l'année.
        J'étudie français parce qu'il aidera avec l'immigration et recherche d'emploi, Québec est une province où tout le monde parle français au lieu de anglais.

    \subsection*{Correction}
        Bonjour, je m'appelle Jz. Je suis étudiant en informatique à l'Université de Waterloo, au Canada. 
        J'ai dix-huit ans et j'habite dans un appartement près du campus.
        J'aime le sport, la musique et la technologie.
        J'adore particuilèrement l'intelligence artificielle et ses applications. 
        Je parle couramment anglais et chinois, et j'apprends le français. 
        Mon objectif est de réussir l'examen à la fin de l'année. J'étudie le français parce qu'il est utile pour l'immigration et pour trouver un emploi au Québec.

\chapter{Day 2}
    \section{Daily Routine}
        En semaine, je me réveille entre six heures en sept heures. Ensuite, je me brosse les dents, me lave ma facile.
        Géneralemment, je mange un pain au fromage et confiture de ceries, et je bois un verre du lait.
        Je me sortis chez moi autour sept heures et je conduis à travail. 
        Je finis travailler à dix-six heures et demi à dix-sept heures et je retourne à ma maison et je me douche juste près ça.
        Je dine géneralement le riz ou la bouille avec un ou deux repas simples.
        Après diner, j'étudie français pour un heure à deux heures en faisant les projets autres personelles.
        Je me couche entre vingt deux heures en midnuit, dépend de que je fais: regarder au vidéo ou jouer à des jeux de vidéo avec mes amis.
        \\
        Sur les weekends, je me réveille autour le même temps comme les jours dans le semaine. J'étudie et je fais des éxercices ou sortir pour visiter mes amis.

    \subsection*{Corrected Version}
        En semaine, je me réveille entre six heures et sept heures.
        Ensuite, je me brosse les dents et je me lave le visage.
        Généralement, je mange du pain avec du formage et de la confiture de cerises, et je bois un verre de lait.
        Je sors de chez moi vers sept heures et je vais au travail.
        Je finis de travailler entre seize heures trente et dix-sept heures, puis je rentre chez moi et je me douche juste après.
        Je dîne généralement du riz ou des bouillies avec un ou deux plats simples. 
        Après le dîner, j'étudie le français pendant une ou deux heures tout en faisant d'autres projets personnels.
        Je me couche entre vingt-deux heures et minuit, selon ce que je fais: regarder des vidéos ou jouer à des jeux vidéo avec mes amis.
        \\
        Le week-end, je me réveille à peu près à la même heure que pendant la semaine. J'étuidie, je fais de l'exercice ou je sors pour rendre visite à mes amis.

    \subsection*{Remarks}
        \begin{itemize}
            \item entre \dots et \dots instead of entre \dots en \dots
            \item se laver le visage, not se laver se visage
            \item sortir de chez moi (irregular conjugation)
            \item aller au travail: je vais au travail en voiture
            \item rentrer chez moi
            \item juste après instead of juste près ça
            \item pendant (for duration, not ``pour'')
            \item personels 
            \item regarder des vidéos (not au)
            \item jeux vidéo (fixed expression)
            \item à peu près à la même heure que (useful structure)
            \item je consacre une à deux heures à l'étude du français (good sentence structure)
            \begin{itemize}
                \item consacrer \dots à \dots (to devote \dots to \dots)
                \item une à deux heures (une expression de durée approximative)
                \item l'étude de quelque chose (the study of something)
            \end{itemize}
            \item jouer: à (for sport and games), de (for musical instruments)
        \end{itemize}

    \subsection*{C1 Version}
        En semaine, je me réveille généralement entre six et sept heures. 
        Après m'être préparé, je prends un petit-déjeuner simple, souvent composé de pain, de fromage et de confiture de cerises, accompagné d'un verre de lait.
        Je quitte ensuite mon appartement vers sept heurs pour aller travailler. 
        En fin d'après-midi, je rentre chez moi, je prends une douche, puis je dîne, le plus souvent de riz ou de bouillies avec quelques plats simples.
        Le soir, je consacre une à deux heures à l'étude du français, tout en avançant sur d'autres projets personnels.
        Je me couche entre vingt-deux heures et minuit, selon mes activités du moment, comme regarder des vidéos ou jouer à des jeux vidéo avec mes amis.
        \\
        Le week-end, mon rythme reste assez similaire, même si je prends aussi le temps de faire de l'exercice et de voir mes amis.

\chapter{Day 3}
    \section{Likes and Dislikes}
        J'adore la musique, je joue du piano, de la saxophone, et j'apprécie écouter de la musique tout le temps.
        J'aime aussi faire du sport, comme jouer au tennis de table, badminton et faire du vélo.
        De plus, j'adore goûter les plats localaux. Je déteste tofu puant et soupe aux haricots.
        Il n'y a pas le plus de chose que je déteste.

    \subsection*{Correction}
        J'adore la musique. Je joue du piano et du saxophone, et j'aime écouter de la musique tout le temps.
        J'aime aussi faire du sport, comme jouer au tennis de table, au badminton et faire du vélo.
        De plus, j'adore goûter des plats locaux. 
        Je déteste le tofu puant et la soupe aux haricots. 
        Il n'y a pas beaucoup de choses que je déteste.

    \subsection*{C1}
        La musique occupe une place très importante dans ma vie. 
        Je joue du piano et du saxophone, et j'aime également écouter de la musique dès que j'en ai l'occasion.
        J'apprécie aussi le sport, en particulier le tennis de table, le badminton et le vélo.
        Par ailleurs, j'aime découvrir des spécialités locales lorsque je peux en goûter. 
        En revanche, certains plats me rebutent, comme le tofu puant ou la soupe aux haricots. 
        De manière généralee, toutefois, il y a peu de choses que je déteste vraiment.

\chapter{Day 4}
    \section{Family and Friends}
        Je n'ai rien des sœurs ou des frères, je suis l'enfant seulment dans ma famille, à cause de loi de seul enfant en Chine.
        Ma mère et mon père ne sont pas jeunes, ils reprendient leur rétraite il y a un couple années.
        Je m'expatrie au Canada en 2017, mes amis de l'école élémentaire sont dans la Chine. 
        Mes amis qui sont dans le Canada, je les rencontre après je m'arrivé au Canada.
        Mmalgré dans le Canada, la plupart des mes amis ont fond chinois, et je les parle dans le choinois le plus souvent plutôt que l'anglais.

    \subsection{Correction}
        Je n'ai ni frère ni sœur. Je suis enfant unique dans ma famille, à cause de la politique de l'enfant unique en Chine.
        Mes parents ne sont plus jeunes; ils ont pris leur retraite il y a quelques années.
        Je suis arrivé au Canada en 2017, donc mes amis de l'école primaire sont en Chine. 
        Quant à mes amis au Canada, je les ai rencontrés après mon arrivée ici. 
        Malgré le fait que je vive au Canada, la plupart de mes amis sont d'origine chinoise,
        et je parle le plus souvent chinois avec eux plutôt qu'anglais.

    \subsection{Remark}
        \begin{itemize}
            \item ni \dots ni \dots
            \item quelques années instead of ``a couple years''
            \item être conjugation for arriver
            \item être d'origine française
            \item parler [language] avec eux or se parler en [language]
            \item si bien que (so thats)
            \item désormais (now)
            \item plutôt que (rather than)
            \item ce qui (which / this fact)
            \item fait que (causes / means / results in the fact that)
            \item ce qui fait que (which means that)
        \end{itemize}

    \subsection{C1}
        Je n'ai ni frère ni sœur, car je suis enfant unique, ce qui s'explique en partie par la politique de l'enfant unique qui a longtemps existé en Chine.
        Mes parents ne sont pas jeunes et ont pris leur retraite il y a déjà quelques années.
        Pour ma part, je suis arrivé au Canada en 2017, si bien que mes amis d'enfance et de l'école primaire vivent toujours en Chine.
        En revanche, les amis que j'ai aujourd'hui au Canada, je les ai recontrés après mon installation ici.
        Même si je vis désormais au Canada, la majorité de mes amis sont d'origine chinoise, ce qui fait que je parle le plus souvent chinois avec eux plutôt qu'anglais.
        
\chapter{Day 5}
    \section{City}
        J'ai habitais en Beijing jusqu'à 2017, l'année je suis arrivé au Canada.
        Beijing est une ville culture. Il y a le palais royal nommé GuGong de la dynastie Qing.
        Maitenant, il est une grand musée avec des plusir expositions uniques et speciaux. 
        Il y a aussi beaucoup des plats et des collations locaux.
        Pour université, je vais à l'université de Waterloo. 
        Le campus est grand et joli, il y a plusir des oies du Canada dans le campus, les errant. 
        Pendant la saison des amours, ils sont devenant trop aggressives, et ils chasseraient les gens qui ses rapproche d'eux.
        Les cours sont un peu lourd dans l'université de Waterloo.

    \subsection{Correection}
        J'ai habité à Beijing jusqu'en 2017, l'année où je suis arrivé au Canada. 
        Beijing est une ville culturelle. Il y a l'ancien palais impérial, appelé la Cité interdite, qui date de la dynastie Qing.
        Maintenant, c'est un grand musée avec plusieurs expositions uniques et spéciales. Il y a aussi beaucoup de plats et de collations locaux.
        Quant à l'université, j'étudie à l'Université de Waterloo. Le campus est grand et joli.
        Il y a aussi plusieurs oies du Canada sur le campus, qui s'y promènent librement. 
        Pendant la saison des amours, elles deviennent trop agressives et elles peuvent poursuivre les gens qui s'approchent d'elles.
        Les cours à l'Université de Waterloo sont un peu lourds.

    \subsection{Remark}
        \begin{itemize}
            \item jusqu'en before a year, sometimes a season and time expressions beginning with en.
        \end{itemize}

    \subsection{C1}
        J'ai vécu à Beijing jusqu'en 2017, année de mon arrivée au Canada. 
        Beijing est une ville particulièrement riche sur le plan culturel et historique.
        On y trouve notamment l'ancien palais impérail, plus connu sous le nom de Cité interdite, qui constitue aujourd'hui un immense musée abritant de nombreuses expositions remarquables.
        La ville est également connue pour la diversité de ses plats et de ses spécialités locales, ce qui la rend très intéressante sur le plan gastronomique.
        Quant à l'Université de Waterloo, où j'étudie actuellement, son campus est vaste et agréable. On y croise souvent des oies du Canada qui se promènent librement.
        Cependant, pendant la saison de reproduction, elles deviennent parfois très agressives et peuvent même poursuivre les personnes qui s'en approchent.
        Enfin, même si j'apprécie beaucoup cette université, je trouve que la charge de travail y est assez lourde.

\chapter{Day 6}
    \section{étudier avec d'autres}
        Je préfére étudier avec d'autres au lieu de étuidier seul. 
        Je pense que la communication est une partie importante pour étudie.
        Cette action non seulement fourni-on une opportunity à apprendre d'autres, mais aussi une opportunité sociale.
        Les autres pouvent donner idées que on ne pense pas. 
        Je apprends toujours dès mes amis pendant que on discute des idées de solutions.
        Quand on se stressé, on peut soulagez la pression en parlant à amis.
        Ainsi, je préfére étudier avec les autres.
    \subsection{Correction}
        Je préfère étuidier avec d'autres au lieu d'étuider seul.
        Je pense que la communication est une partie importante de l'apprentissage.
        Cette façon d'étudier fournit non seulement une occasion d'apprendre des autres,
        mais aussi une opportunité de socialiser.
        Les autres peuvent donner des idées auxquelles on ne pense pas.
        J'apprends toujours de mes amis pendant que nous discutons des solutions.
        Quand on est stressé, on peut soulager la pression en parlant à ses amis.
        Ainsi, je préfère étudier avec les autres.

    \subsection{C1}
        Je préfère étuider avec d'autres plutôt que seul.
        À mon avis, la communication joue un rôle essentiel dans l'apprentissage.
        D'abord, le travail en groupe permet non seulement d'apprendre des autres, 
        mais aussi d'échanger des points de vue que l'on n'aurait pas envisagés seul.
        Ensuite, il offre aussi une dimension sociale, ce qui peut rendre l'étude plus motivante et moins stressante.
        Lorsque l'on rencontre des difficultés, le fait de parler avec ses amis aide souvent à clarifier ses idées et à relâcher la pression.
        En conclusion, même si l'étude individuelle a aussi ses avantages, je trouve qu'étudier avec d'autres est souvent plus efficace et plus agréable.





\end{document}